% Educational mapping, not a reproduction of the original paper figure.
% Center and each arrow are textual; color is not the sole communication cue.
\tagstructbegin{tag=Figure,alttext={左右两路的 Encoder--Decoder Transformer 示意，左侧 Encoder 与右侧 Cross-Attention 用灰色虚线标示在 Decoder-only 中不保留：左路源 Token 经 Embedding、双向 Encoder 自注意力、FFN、残差和归一化成为 Encoder 输出，供交叉注意力读取为 Key 与 Value；右路目标前缀先经过因果遮挡的 Decoder 自注意力，再以 Decoder 状态为 Query 对左路 Encoder 输出做 Cross-Attention，随后经过残差、归一化、FFN、输出头和 Softmax。图中框还标明教材章节；Decoder-only 模型去掉独立 Encoder 与交叉注意力，但保留因果自注意力、FFN 与输出头。}}
\begin{center}
\begin{tikzpicture}[>=Latex,
 mapbox/.style={draw=linegray,fill=paper,rounded corners=2pt,minimum width=52mm,
    minimum height=12mm,text width=46mm,align=center,font=\scriptsize},
 legacy/.style={mapbox,draw=muted,dashed,fill=linegray!22,text=muted},
 link/.style={->,line width=.7pt}]
\node[legacy] (src) at (-3.35,2.55) {源 Token / Embedding / 位置\\第 12、14、21 章};
\node[legacy] (enc) at (-3.35,1.05) {Encoder Self-Attention + FFN\\残差／归一化，双向源输入\\第 19、20、22、23 章};
\node[legacy] (memory) at (-3.35,-.45) {Encoder 输出：源的 Key / Value};
\node[mapbox] (tgt) at (3.35,2.55) {目标前缀 Token / Embedding / 位置\\第 12、14、21 章};
\node[mapbox] (mask) at (3.35,1.05) {Masked Decoder Self-Attention\\因果可见\\第 19--22 章};
\node[legacy] (cross) at (3.35,-.45) {Cross-Attention\\Query 取自 Decoder，K/V 取自 Encoder};
\node[mapbox] (ffn) at (3.35,-1.95) {Add / Norm 与逐位置 FFN\\第 22、23 章};
\node[mapbox] (soft) at (3.35,-3.45) {输出头、logits、Softmax\\第 23--26 章};
\draw[link,muted,dashed] (src)--(enc);
\draw[link,muted,dashed] (enc)--(memory);
\draw[link,blue] (tgt)--(mask);
\draw[link,muted,dashed] (mask)--(cross);
% Blue bypass denotes the simplified Decoder-only route without Cross-Attention.
\draw[link,blue] (mask.east)--++(0.65,0) |- node[pos=.72,right,font=\scriptsize,text=blue]{保留路径} (ffn.east);
\draw[link,muted,dashed] (cross)--(ffn);
\draw[link,blue] (ffn)--(soft);
\draw[link,muted,dashed] (memory.east)--node[above,font=\scriptsize]{源 K/V} (cross.west);
\end{tikzpicture}
\par\smallskip{\sffamily\footnotesize\color{muted}灰色虚线：Decoder-only 中移除的独立 Encoder 与 Cross-Attention；蓝色实线：保留的主路径。}
\end{center}
\begin{center}
\captionof{figure}{原论文的 Encoder--Decoder 路径对照：Decoder-only 主线保留因果自注意力、FFN 和输出头，灰色虚线表示从 Decoder-only 中移除的 Encoder 与 Cross-Attention}
\label{fig:original-transformer-map}
\end{center}
\tagstructend
